\documentclass[11pt]{article}

% ── packages ──────────────────────────────────────────────────────────────────
\usepackage[margin=1in]{geometry}
\usepackage{amsmath, amssymb}
\usepackage{graphicx}
\usepackage{booktabs}
\usepackage{hyperref}
\usepackage{cleveref}
\usepackage{caption}
\usepackage{subcaption}
\usepackage{parskip}
\usepackage{microtype}
\usepackage{xcolor}
\usepackage{listings}

% ── listings style ────────────────────────────────────────────────────────────
\lstset{
  basicstyle=\ttfamily\small,
  keywordstyle=\color{blue},
  commentstyle=\color{gray},
  breaklines=true,
  frame=single,
  language=Python,
}

% ── title ─────────────────────────────────────────────────────────────────────
\title{\textbf{Probability of Collision in Low Earth Orbit}\\[0.4em]
       \large Synthetic Scenario Generation, TCA Finding,\\
       and the Fowler (1993) Analytic Method}
\author{Grace Kim}
\date{April 2026}

\begin{document}
\maketitle
\tableofcontents
\newpage

% ══════════════════════════════════════════════════════════════════════════════
\section{Introduction}
\label{sec:intro}
% ══════════════════════════════════════════════════════════════════════════════

The probability of collision ($P_c$) between two orbiting objects is a central
quantity in space situational awareness (SSA).  Given the growing number of
active satellites and debris objects in low Earth orbit (LEO), fast and accurate
$P_c$ calculations are critical for conjunction screening and collision avoidance
maneuver decisions.

This report documents the first phase of a systematic study of $P_c$ methods.
The goals are threefold:
\begin{enumerate}
  \item Build a \textbf{synthetic conjunction generator} that produces
        realistic LEO encounter scenarios with known ground truth, enabling
        rigorous testing of any $P_c$ method.
  \item Implement a \textbf{TCA finder} that locates the Time of Closest
        Approach from initial conditions, as required by all analytic $P_c$
        methods.
  \item Implement the \textbf{Fowler (1993)} analytic $P_c$ method and
        characterize its behavior across encounter geometries and covariance
        regimes.
\end{enumerate}

All code is written in Python, using the \texttt{brahe} astrodynamics library
for orbit propagation and coordinate frame transformations.

% ══════════════════════════════════════════════════════════════════════════════
\section{Background}
\label{sec:background}
% ══════════════════════════════════════════════════════════════════════════════

\subsection{Reference Frames}

Two frames are used throughout.

\paragraph{Earth-Centered Inertial (ECI).}  Origin at Earth's center, axes
fixed to the celestial sphere (J2000 epoch).  The full spacecraft state is the
6-vector $\mathbf{x} = [\mathbf{r}^\top, \mathbf{v}^\top]^\top$ in meters and
meters per second.

\paragraph{Radial-Tangential-Normal (RTN).}  A body-fixed frame centered on a
reference spacecraft (SC1).  The axes are:
\begin{align}
  \hat{\mathbf{r}} &= \frac{\mathbf{r}_1}{|\mathbf{r}_1|}
    && \text{(radial — toward SC1 from Earth center)}, \\
  \hat{\mathbf{n}} &= \frac{\mathbf{r}_1 \times \mathbf{v}_1}
                           {|\mathbf{r}_1 \times \mathbf{v}_1|}
    && \text{(normal — orbit angular momentum direction)}, \\
  \hat{\mathbf{t}} &= \hat{\mathbf{n}} \times \hat{\mathbf{r}}
    && \text{(tangential/along-track)}.
\end{align}
RTN is natural for expressing inter-spacecraft relative states and for
parameterizing covariance matrices because tracking errors are largest along
track.

\subsection{Keplerian Orbital Elements}

SC1's initial orbit is specified by six classical elements
$(a, e, i, \Omega, \omega, M_0)$:
semi-major axis, eccentricity, inclination, right ascension of the ascending
node, argument of perigee, and mean anomaly at epoch.  All angles are in
degrees; $a$ is in meters.  The \texttt{brahe} function
\texttt{state\_koe\_to\_eci} converts these to an ECI state vector.

\subsection{Conjunction Geometry}

A \emph{conjunction} is a close approach event between two spacecraft.  The key
scalar quantities at the Time of Closest Approach (TCA) are:
\begin{itemize}
  \item \textbf{Miss distance} $d = |\mathbf{r}_1 - \mathbf{r}_2|_{\text{TCA}}$
        (meters).
  \item \textbf{Relative speed} $v_\text{rel} = |\mathbf{v}_1 - \mathbf{v}_2|_{\text{TCA}}$
        (m/s).
  \item \textbf{Hard-body radius (HBR)} $r_\text{hb}$: the combined physical
        cross-section radius of the two objects (meters).  A collision is said
        to occur when $d \le r_\text{hb}$.
\end{itemize}

% ══════════════════════════════════════════════════════════════════════════════
\section{Synthetic Conjunction Generation}
\label{sec:conjunction}
% ══════════════════════════════════════════════════════════════════════════════

\subsection{Design Philosophy}

Rather than choosing initial conditions and hoping a close approach occurs
within a search window, the generator works \emph{backwards from TCA}:

\begin{enumerate}
  \item Define SC1's Keplerian orbit at TCA.
  \item Place SC2 relative to SC1 at TCA using the desired miss distance
        $d$ and relative speed $v_\text{rel}$, expressed in the RTN frame.
        This guarantees the encounter geometry by construction.
  \item Propagate both spacecraft \emph{backwards} 24 hours from TCA using
        a numerical integrator to obtain $T{=}0$ initial conditions.
\end{enumerate}

Because the TCA geometry is designed explicitly, the ground truth
(exact miss distance, exact TCA epoch, and TCA states) is always known.  This
is essential for validating TCA-finding algorithms and $P_c$ methods.

\subsection{SC1 Default Orbit}

SC1 is placed in a representative LEO orbit at an altitude of 550\,km:

\begin{center}
\begin{tabular}{ll}
  \toprule
  Element & Value \\
  \midrule
  Semi-major axis $a$ & $R_\oplus + 550{,}000$\,m \\
  Eccentricity $e$ & 0.001 \\
  Inclination $i$ & 55° \\
  RAAN $\Omega$ & 20° \\
  Arg.\ of perigee $\omega$ & 0° \\
  Mean anomaly $M_0$ & 0° (at TCA) \\
  \bottomrule
\end{tabular}
\end{center}

\subsection{SC2 Placement in RTN}

SC2's relative state at TCA is constructed so that:
\begin{align}
  |\mathbf{r}_\text{rel}|_{\text{TCA}} &= d  && \text{(desired miss distance)}, \\
  |\mathbf{v}_\text{rel}|_{\text{TCA}} &= v_\text{rel} && \text{(desired relative speed)}, \\
  \mathbf{r}_\text{rel} \cdot \mathbf{v}_\text{rel} &= 0
  && \text{(TCA condition — relative motion perpendicular to miss vector)}.
\end{align}

The perpendicularity condition is enforced by projecting the velocity out of the
miss-vector direction:
\begin{equation}
  \mathbf{v}_\perp = \mathbf{v}_\text{nom} -
    \frac{\mathbf{v}_\text{nom} \cdot \mathbf{r}_\text{rel}}
         {|\mathbf{r}_\text{rel}|^2}\,\mathbf{r}_\text{rel},
  \qquad
  \mathbf{v}_\text{rel} = v_\text{rel}\,\frac{\mathbf{v}_\perp}{|\mathbf{v}_\perp|}.
\end{equation}

\subsection{Conjunction Types}
\label{sec:ctypes}

Four encounter geometries are implemented.  \Cref{tab:ctypes} summarizes the
RTN placement and physical interpretation of each.  Three of them
(crossing, head-on, overtaking) represent distinct orbital mechanics regimes.
The fourth — near-miss — uses the \emph{same} crossing geometry but with a
10\,m miss distance (equal to the HBR), making it a stress test for $P_c$
methods rather than a distinct encounter type.

\begin{table}[h]
\centering
\caption{Conjunction types, RTN geometry at TCA, and physical interpretation.}
\label{tab:ctypes}
\begin{tabular}{lllp{5.8cm}}
  \toprule
  Type & Miss direction & $\mathbf{v}_\text{rel}$ dominant axis & Physical meaning \\
  \midrule
  Crossing  & $+\hat{\mathbf{t}}$ & $-\hat{\mathbf{n}}$ &
    Two orbits in different planes; $v_\text{rel}$ is cross-track.
    Most common real-world conjunction type.  Encounter plane is $\approx$R-T,
    so the large along-track uncertainty projects in $\Rightarrow$ low $P_c$. \\[4pt]
  Head-on   & $+\hat{\mathbf{t}}$ & retrograde ($\approx -\hat{\mathbf{t}}$) &
    SC2 in a retrograde orbit; $v_\text{rel} \approx 2 \times v_\text{orb}
    \approx 15$\,km/s.  Encounter plane is $\approx$R-N;
    along-track uncertainty is \emph{excluded} $\Rightarrow$ compact
    projected covariance, higher $P_c$. \\[4pt]
  Overtaking & $+\hat{\mathbf{t}}$ & $-\hat{\mathbf{t}}$ &
    SC2 slightly ahead in the same orbital plane; SC1 slowly catching up;
    $v_\text{rel}$ is small ($\sim$50\,m/s).  Encounter plane is also $\approx$R-N,
    but the large miss distance (1000\,m) keeps $P_c$ low. \\[4pt]
  Near-miss  & $+\hat{\mathbf{t}}$ & $-\hat{\mathbf{n}}$ &
    Identical geometry to crossing but miss distance $= 10$\,m $\approx$ HBR.
    Same encounter plane, same wide projected covariance — but the Gaussian
    mean sits \emph{at} the HBR boundary, giving $P_c \sim 10^{-3}$. \\
  \bottomrule
\end{tabular}
\end{table}

\paragraph{Why near-miss is distinct from crossing despite identical geometry.}
The crossing and near-miss scenarios use the same conjunction type code path and
the same relative velocity direction.  The difference is entirely in the miss
distance: 500\,m vs.\ 10\,m.  In the encounter-plane view, both have the same
elongated covariance ellipse (wide along the T axis).  For the 500\,m crossing,
the Gaussian center is far from the HBR disk — the tail barely reaches it, so
$P_c \sim 10^{-14}$.  For the near-miss, the center sits right at the disk edge,
so the integral captures a meaningful fraction of the Gaussian and $P_c \sim 10^{-3}$.
This pair isolates the effect of miss distance from geometry (compare
\Cref{fig:conj_crossing,fig:conj_nearmiss}).

\paragraph{Special case: head-on.}  The RTN frame subtends only $\sim$63° between
along-track and cross-track; a true 180° retrograde orbit cannot be expressed
as a small RTN offset.  Instead, \texttt{\_sc2\_eci\_head\_on} constructs SC2
directly in ECI by decomposing SC1's velocity into R, T, N components and
negating the T component, producing an orbit with the same altitude but opposite
direction:
\begin{equation}
  \mathbf{v}_2 = \mathbf{v}_{1,R} - \mathbf{v}_{1,T} + \mathbf{v}_{1,N}.
\end{equation}

\subsection{Per-Scenario Diagnostics}
\label{sec:vis_4panel}

\Cref{fig:conj_crossing,fig:conj_headon,fig:conj_overtaking,fig:conj_nearmiss}
show the four-panel diagnostic plot for each scenario.  The panels are:
\begin{enumerate}
  \item \textbf{RTN trajectory} — $\Delta R$, $\Delta T$, $\Delta N$ vs.\ time.
        The axis that crosses zero near $t=24$\,h identifies the approach
        direction and the conjunction type.
  \item \textbf{Miss distance vs.\ time} — scalar separation with a vertical
        line at the TCA found by \texttt{find\_tca} and a horizontal line at
        the HBR.  The width of the minimum shows how sharply the spacecraft
        close in.
  \item \textbf{Encounter-plane projection} — 1$\sigma$/2$\sigma$/3$\sigma$
        covariance ellipses, the HBR disk (origin), and the miss point (cross).
        The shape of the ellipse is the key quantity: compact for head-on
        (high $P_c$), elongated for crossing/overtaking (low $P_c$).
        When the miss distance is large relative to the ellipse spread (crossing,
        overtaking), a zoomed inset shows the HBR disk at the origin.
  \item \textbf{RTN covariance cross-sections} — position uncertainty
        ellipses in the R-T, R-N, and T-N planes using the default 1-$\sigma$
        values.  The T axis dominates, as expected for LEO tracking errors.
\end{enumerate}

\begin{figure}[htb]
  \centering
  \includegraphics[width=\textwidth]{../plots/conjunctions_crossing.png}
  \caption{\textbf{Crossing conjunction} (500\,m miss, $v_\text{rel} = 7{,}000$\,m/s,
    N-dominant approach).
    \emph{RTN trajectory (panel 1):} the N component oscillates and crosses zero
    at TCA (red dashed line), while T and R components show the long 24-hour
    approach arc.
    \emph{Miss distance (panel 2):} many oscillations over 24 hours with a
    smooth final dip — characteristic of a crossing orbit whose ground tracks
    approach each other gradually.
    \emph{Encounter plane (panel 3):} the covariance ellipse is elongated
    horizontally (along the T axis) because $\mathbf{v}_\text{rel} \approx \hat{\mathbf{n}}$
    places the encounter plane $\approx$R-T, projecting the full 500\,m
    along-track sigma into the 2D view.  The miss point (X) is far from the
    10\,m HBR disk (zoomed inset, bottom-left).  Result: $P_c \sim 10^{-14}$.}
  \label{fig:conj_crossing}
\end{figure}

\begin{figure}[htb]
  \centering
  \includegraphics[width=\textwidth]{../plots/conjunctions_head_on.png}
  \caption{\textbf{Head-on conjunction} (200\,m miss, $v_\text{rel} \approx 15$\,km/s,
    retrograde SC2).
    \emph{RTN trajectory (panel 1):} T and R components oscillate over 24 hours
    with high amplitude (retrograde orbit sweeps large RTN excursions); N is small.
    \emph{Miss distance (panel 2):} rapid oscillations reflecting the high-speed
    retrograde geometry, with a sharp minimum at TCA — the miss distance profile
    collapses quickly because $v_\text{rel} \approx 15$\,km/s.
    \emph{Encounter plane (panel 3):} the covariance ellipse is compact and
    nearly circular.  With $\mathbf{v}_\text{rel} \approx -\hat{\mathbf{t}}$,
    the encounter plane is $\approx$R-N, projecting only the small R (100\,m)
    and N (50\,m) sigmas into the 2D view.  The full 500\,m along-track sigma
    is perpendicular to the plane and excluded.  The HBR disk is visible at scale,
    with the miss point close enough to contribute significant probability.
    Result: $P_c \sim 5\times10^{-4}$.}
  \label{fig:conj_headon}
\end{figure}

\begin{figure}[htb]
  \centering
  \includegraphics[width=\textwidth]{../plots/conjunctions_overtaking.png}
  \caption{\textbf{Overtaking conjunction} (1000\,m miss, $v_\text{rel} = 50$\,m/s,
    T-dominant approach).
    \emph{RTN trajectory (panel 1):} both R and T components slowly decline toward
    zero as SC1 catches SC2.  The oscillation pattern is slow, reflecting the
    small relative speed.
    \emph{Miss distance (panel 2):} many smooth oscillations over 24 hours as the
    two spacecraft slowly converge, with a broad minimum — a slow, gentle approach.
    \emph{Encounter plane (panel 3):} despite $\mathbf{v}_\text{rel} \approx -\hat{\mathbf{t}}$
    pointing the encounter plane toward R-N (same as head-on), the miss distance
    is 1000\,m.  The ellipse is elongated because the overtaking scenario has
    the same default covariance, but with $v_\text{rel}$ somewhat off-axis from
    pure along-track at this seed, some T-axis uncertainty projects in.
    The zoomed inset shows the 10\,m HBR disk at the origin — far from the
    1000\,m miss point — explaining the low $P_c$.}
  \label{fig:conj_overtaking}
\end{figure}

\begin{figure}[htb]
  \centering
  \includegraphics[width=\textwidth]{../plots/conjunctions_near_miss.png}
  \caption{\textbf{Near-miss conjunction} (10\,m miss, $v_\text{rel} = 500$\,m/s,
    crossing geometry — same as \Cref{fig:conj_crossing} except for miss distance).
    \emph{RTN trajectory (panel 1):} identical qualitative shape to the crossing
    scenario.  The N component approaches zero at TCA; only the final miss distance
    differs.
    \emph{Miss distance (panel 2):} smooth approach to 10\,m at TCA — the HBR
    line (horizontal dashes) is nearly touching the minimum.
    \emph{Encounter plane (panel 3):} the same elongated covariance ellipse as
    the crossing case, but now the miss point (X) sits nearly at the edge of
    the HBR disk.  Even though the ellipse is wide (the 500\,m along-track sigma
    still projects in), the center is close enough to the disk that a non-trivial
    fraction of the Gaussian mass falls inside.  Result: $P_c \sim 10^{-3}$,
    demonstrating that miss distance alone can drive $P_c$ up by 11 orders of
    magnitude relative to the geometrically identical crossing scenario.}
  \label{fig:conj_nearmiss}
\end{figure}

\clearpage

\subsection{3D RTN Relative Motion}
\label{sec:vis_3d}

\Cref{fig:3d_overview} shows the 3D relative trajectory of SC2 with SC1 fixed
at the origin in the RTN frame.  Color varies from light red ($T{=}0$, 24\,h
before TCA) to dark red (TCA).

\paragraph{Why the crossing trajectory looks like a spiral.}
The crossing scenario shows SC2 tracing a large, slowly-shrinking spiral in
the RTN frame (\Cref{fig:3d_overview}, top-left).  This is not SC2 in a
physical orbit around SC1 — it is an artifact of the RTN frame rotating with
SC1's orbital motion.  SC1 completes roughly 14 orbits in 24 hours at 550\,km
altitude.  Each orbit, the RTN frame rotates once with respect to inertial space.
SC2 is in a fixed (non-rotating) orbit at a slightly different inclination, so
its relative position in the RTN frame traces an ellipse once per orbit.
Over 24 orbits the ellipses shift slightly as the crossing point slowly
approaches, producing the spiral appearance.  The actual physical motion is two
satellites in fixed orbits whose paths happen to cross once — the spiral is a
frame effect, not a real trajectory.

\begin{figure}[htb]
  \centering
  \includegraphics[width=\linewidth]{../plots/conjunctions_3d_orbits.png}
  \caption{3D RTN relative trajectories for all four conjunction types.
           SC1 is at the origin; SC2's path is colored light$\to$dark red
           as time progresses toward TCA.
           \textbf{Crossing (top-left):} SC2 appears to spiral inward over
           24 hours — a frame artifact (see text); the actual close approach
           is the tiny dark-red cluster near the origin.
           \textbf{Head-on (top-right):} a tight ellipse reflecting the
           retrograde orbit geometry; the approach and departure are rapid
           ($\sim$15\,km/s) so the TCA cluster is compact.
           \textbf{Overtaking (bottom-left):} SC2 slides slowly along $+T$
           toward SC1, then begins to lap SC1 as the relative speed is low
           ($\sim$50\,m/s); the 24-hour arc is a long, nearly-linear tube.
           \textbf{Near-miss (bottom-right):} same spiral topology as crossing
           but with a 10\,m miss distance; the TCA cluster collapses to a point
           essentially at the origin.}
  \label{fig:3d_overview}
\end{figure}

% ══════════════════════════════════════════════════════════════════════════════
\section{Finding the Time of Closest Approach}
\label{sec:tca}
% ══════════════════════════════════════════════════════════════════════════════

\subsection{Why TCA Must Be Found Numerically}

All analytic $P_c$ methods — including Fowler (1993) — require the spacecraft
states \emph{at TCA}: the relative position $\mathbf{r}_\text{rel}$ and relative
velocity $\mathbf{v}_\text{rel}$ at the moment of closest approach.  These are
needed to define the encounter plane and compute the miss vector $\mathbf{m}$.

In practice, however, the inputs to a conjunction screening system are not states
at TCA.  They are states at some catalog epoch — call it $T{=}0$ — which is
typically 12--48 hours before the anticipated close approach.  You do not know in
advance \emph{when} closest approach occurs; that is what you are trying to find.

There is no closed-form expression for the TCA of two objects in general orbits.
Even for two-body (Keplerian) motion, the relative trajectory is a complicated
function of the six initial orbital elements of each object, and the time of
minimum separation does not have an analytic solution.  It must be found
numerically by propagating both trajectories and searching for the minimum
of $d(t) = |\mathbf{r}_1(t) - \mathbf{r}_2(t)|$.

This is non-trivial because:
\begin{itemize}
  \item Each evaluation of $d(t)$ requires integrating two sets of differential
        equations from $T{=}0$ to time $t$ — this is expensive.
  \item The function $d(t)$ can have multiple local minima over a 24-hour window
        (e.g.\ for slow overtaking scenarios where the two satellites are nearly
        co-orbital and approach each other repeatedly).
  \item Standard gradient-based minimizers fail because the numerical integrator
        produces slightly noisy function values.
\end{itemize}

A robust two-stage algorithm — coarse grid search followed by scalar minimization
— handles all of these challenges.

\subsection{Problem Statement}

Given initial conditions $(\mathbf{x}_1(0), \mathbf{x}_2(0))$ at $T{=}0$,
find the time $t^*$ in the planning horizon $[0,\; T_\text{horizon}]$ that
minimizes the miss distance:
\begin{equation}
  t^* = \arg\min_{t \in [0,\;T_\text{horizon}]}
    \bigl|\mathbf{r}_1(t) - \mathbf{r}_2(t)\bigr|,
\end{equation}
where $\mathbf{r}_i(t)$ is obtained by numerical integration of the equations
of motion from $T{=}0$.

\subsection{Algorithm}

The search proceeds in two stages.

\paragraph{Stage 1 — Coarse grid.}
A single pair of propagators is advanced sequentially through $N_c = 200$
evenly-spaced time nodes spanning the 24-hour window.  At each node $t_k$,
the miss distance
\begin{equation}
  d_k = |\mathbf{r}_1(t_k) - \mathbf{r}_2(t_k)|
\end{equation}
is recorded.  Using a single pair of propagators (rather than creating fresh
ones at each node) means only $N_c$ integration steps are performed in total —
$O(N_c)$ work.  The best node $k^*$ is identified as $\arg\min_k d_k$.

\paragraph{Stage 2 — Fine minimization.}
Brent's method (\texttt{scipy.optimize.minimize\_scalar} with
\texttt{method='bounded'}) refines the minimum within the interval
$[t_{k^*} - \Delta t,\; t_{k^*} + \Delta t]$ where $\Delta t$ is one coarse
grid step ($\approx 7.2$\,min).  The search is anchored at the state
$({\mathbf{x}}_1(t_{k^*}), \mathbf{x}_2(t_{k^*}))$ so the integrator only
propagates a short arc, not the full 24-hour window.  The tolerance is set to
$\epsilon = 1$\,s.

\paragraph{Typical performance.}
A single \texttt{find\_tca} call takes 3--4\,s on a modern laptop: the coarse
grid performs 200 propagation steps; the fine minimization performs
$O(10)$ short-arc propagations.

\subsection{TCA Timing and Miss Distance Accuracy}

\Cref{tab:tca_accuracy} compares the TCA found by the algorithm against the
ground truth planted by the conjunction generator.

\begin{table}[h]
\centering
\caption{TCA finder accuracy across all scenario types.}
\label{tab:tca_accuracy}
\begin{tabular}{lrrrr}
  \toprule
  Scenario & Miss (m) & $v_\text{rel}$ (m/s) & Timing error (s) & Miss error (m) \\
  \midrule
  Crossing   & 500 & 7000 & $<5$ & $<5$ \\
  Head-on    & 200 & 15000 & $<5$ & $<5$ \\
  Overtaking & 1000 & 15 & $<5$ & $<5$ \\
  Near-miss  & 10 & 500 & $<5$ & $<5$ \\
  \bottomrule
\end{tabular}
\end{table}

% ══════════════════════════════════════════════════════════════════════════════
\section{Covariance Model}
\label{sec:covariance}
% ══════════════════════════════════════════════════════════════════════════════

\subsection{Covariance Representation}

Each spacecraft's state uncertainty is represented as a $6\times6$
positive-definite covariance matrix $\mathbf{C}$ in ECI, with units
$[\text{m}^2,\;\text{m}^2/\text{s}^2,\;\text{m}\cdot\text{m/s}]$ for the
position-position, velocity-velocity, and cross blocks respectively.

\subsection{RTN Parameterization}

Tracking errors are naturally anisotropic in the RTN frame: along-track
uncertainty dominates because orbit determination error grows fastest in the
direction of motion.  The covariance is therefore specified as diagonal in RTN
and then rotated to ECI:

\begin{enumerate}
  \item Construct a diagonal $6\times6$ RTN covariance from six 1-$\sigma$
        values $(\sigma_R, \sigma_T, \sigma_N, \sigma_{\dot R}, \sigma_{\dot T},
        \sigma_{\dot N})$:
        \begin{equation}
          \mathbf{C}_\text{RTN} = \mathrm{diag}
            \bigl(\sigma_R^2,\; \sigma_T^2,\; \sigma_N^2,\;
                  \sigma_{\dot R}^2,\; \sigma_{\dot T}^2,\; \sigma_{\dot N}^2\bigr).
        \end{equation}
  \item Fetch the $3\times3$ rotation matrix
        $\mathbf{R} = \texttt{brahe.rotation\_rtn\_to\_eci}(\mathbf{x}_\text{ECI})$,
        which maps RTN unit vectors to ECI columns.
  \item Build a $6\times6$ block-diagonal rotation:
        \begin{equation}
          \mathbf{T} = \begin{pmatrix} \mathbf{R} & \mathbf{0} \\
                                      \mathbf{0} & \mathbf{R} \end{pmatrix}.
        \end{equation}
  \item Rotate: $\mathbf{C}_\text{ECI} = \mathbf{T}\,\mathbf{C}_\text{RTN}\,\mathbf{T}^\top$.
\end{enumerate}

The trace of the position block is invariant under this rotation:
$\mathrm{tr}(\mathbf{C}_\text{ECI}[:3,:3]) = \sigma_R^2 + \sigma_T^2 + \sigma_N^2$.

\subsection{Default 1-$\sigma$ Values}

The default uncertainty profile is representative of well-tracked LEO objects:

\begin{center}
\begin{tabular}{lll}
  \toprule
  Component & Position (m) & Velocity (m/s) \\
  \midrule
  Radial (R) & 100 & 0.10 \\
  Along-track (T) & 500 & 0.50 \\
  Cross-track (N) & 50 & 0.05 \\
  \bottomrule
\end{tabular}
\end{center}

\subsection{Simplification: Covariance Constructed at TCA}

In a real conjunction pipeline, the covariance at $T{=}0$ would be propagated
forward to TCA using a state transition matrix (STM):
\begin{equation}
  \mathbf{C}_\text{TCA} = \boldsymbol{\Phi}(t_\text{TCA}, 0)\;\mathbf{C}_0\;
  \boldsymbol{\Phi}(t_\text{TCA}, 0)^\top,
\end{equation}
where $\boldsymbol{\Phi}$ is the $6\times6$ Jacobian of the propagated state
with respect to the initial state.  Over 24 hours, along-track uncertainty can
inflate by a factor of 10--100 due to differential drag and gravitational
perturbations.

For the purposes of this study (testing $P_c$ methods in isolation), covariances
are constructed directly at TCA from fixed RTN 1-$\sigma$ assumptions.  This is
correct for evaluating the Fowler algorithm independent of orbital mechanics
errors, but a full end-to-end pipeline must propagate the covariance.  STM
propagation is planned as a future addition (see \S\ref{sec:future}).

% ══════════════════════════════════════════════════════════════════════════════
\section{Probability of Collision: Fowler (1993) Method}
\label{sec:fowler}
% ══════════════════════════════════════════════════════════════════════════════

\subsection{Overview}

The Fowler (1993) method reduces the 3D collision problem to a 2D integral by
projecting all relevant quantities onto the \emph{encounter plane} — the plane
perpendicular to the relative velocity vector at TCA.  The key physical insight
is that, during the brief encounter, the relative motion is dominated by the
(nearly constant) relative velocity; motion within the encounter plane can be
treated as instantaneous relative to the collision timescale.

\subsection{Encounter Plane Construction}

Given ECI states $\mathbf{x}_1, \mathbf{x}_2$ at TCA, define:
\begin{align}
  \mathbf{r}_\text{rel} &= \mathbf{r}_1 - \mathbf{r}_2, \\
  \mathbf{v}_\text{rel} &= \mathbf{v}_1 - \mathbf{v}_2.
\end{align}

The encounter-plane orthonormal basis is:
\begin{align}
  \hat{\mathbf{z}} &= \frac{\mathbf{v}_\text{rel}}{|\mathbf{v}_\text{rel}|}
    && \text{(normal to encounter plane, along relative velocity)}, \\
  \hat{\mathbf{x}} &= \frac{\mathbf{r}_\text{rel} -
      (\mathbf{r}_\text{rel} \cdot \hat{\mathbf{z}})\hat{\mathbf{z}}}
    {\bigl|\mathbf{r}_\text{rel} - (\mathbf{r}_\text{rel} \cdot \hat{\mathbf{z}})\hat{\mathbf{z}}\bigr|}
    && \text{(in-plane, toward miss vector)}, \\
  \hat{\mathbf{y}} &= \hat{\mathbf{z}} \times \hat{\mathbf{x}}
    && \text{(completes right-hand frame)}.
\end{align}

If $\mathbf{r}_\text{rel} \parallel \mathbf{v}_\text{rel}$ (zero impact parameter),
$\hat{\mathbf{x}}$ is chosen as an arbitrary unit vector perpendicular to
$\hat{\mathbf{z}}$.

\subsection{2D Projection}

Let $\mathbf{B} = [\hat{\mathbf{x}},\, \hat{\mathbf{y}}]^\top \in \mathbb{R}^{2\times3}$.
The combined 3D position covariance is
\begin{equation}
  \mathbf{C}_\text{pos} = \mathbf{C}_1[{:}3,{:}3] + \mathbf{C}_2[{:}3,{:}3].
\end{equation}
Projecting onto the encounter plane:
\begin{align}
  \mathbf{C}_\text{2D} &= \mathbf{B}\,\mathbf{C}_\text{pos}\,\mathbf{B}^\top
    \in \mathbb{R}^{2\times2}, \label{eq:c2d}\\
  \mathbf{m} &= \mathbf{B}\,\mathbf{r}_\text{rel} \in \mathbb{R}^2.
    \label{eq:miss2d}
\end{align}

\subsection{The $P_c$ Integral}

The probability of collision is the integral of a 2D Gaussian centered at
$\mathbf{m}$ with covariance $\mathbf{C}_\text{2D}$ over the hard-body disk
of radius $r_\text{hb}$:

\begin{equation}
  \boxed{
  P_c = \int\!\!\int_{x^2 + y^2 \le r_\text{hb}^2}
    \mathcal{N}\!\left([x,\,y]^\top;\;\mathbf{m},\;\mathbf{C}_\text{2D}\right)
    \, dx\, dy.}
  \label{eq:fowler_pc}
\end{equation}

Here $\mathcal{N}(\mathbf{u};\;\boldsymbol{\mu},\;\boldsymbol{\Sigma})$ denotes
the 2D Gaussian PDF evaluated at $\mathbf{u}$:
\begin{equation}
  \mathcal{N}(\mathbf{u};\;\boldsymbol{\mu},\;\boldsymbol{\Sigma})
  = \frac{1}{2\pi\sqrt{|\boldsymbol{\Sigma}|}}
    \exp\!\left(-\tfrac{1}{2}
    (\mathbf{u}-\boldsymbol{\mu})^\top\boldsymbol{\Sigma}^{-1}
    (\mathbf{u}-\boldsymbol{\mu})\right).
\end{equation}

The integral \eqref{eq:fowler_pc} is evaluated numerically using
\texttt{scipy.integrate.dblquad}, parameterized as:
\begin{equation}
  P_c = \int_{-r_\text{hb}}^{r_\text{hb}}
    \int_{-\sqrt{r_\text{hb}^2 - y^2}}^{+\sqrt{r_\text{hb}^2 - y^2}}
    \mathcal{N}([x,y]^\top;\;\mathbf{m},\;\mathbf{C}_\text{2D})\;dx\;dy.
  \label{eq:dblquad}
\end{equation}

Absolute and relative tolerances are set to $10^{-10}$ and $10^{-8}$,
respectively.  The result is clamped to $[0, 1]$ to guard against tiny
numerical overshoots.

\subsection{Why the Encounter Geometry Dominates $P_c$}

The encounter plane construction makes explicit why two encounters with the
same miss distance and same covariance can have radically different $P_c$:
\emph{it is the direction of $\mathbf{v}_\text{rel}$ that determines which
axes of the covariance project into $\mathbf{C}_\text{2D}$.}

\begin{itemize}
  \item \textbf{Head-on} ($\mathbf{v}_\text{rel} \approx \hat{\mathbf{t}}$):
    the encounter plane is roughly the R-N plane.  The large along-track
    uncertainty ($\sigma_T = 500$\,m) is \emph{excluded} from $\mathbf{C}_\text{2D}$
    because it lies along $\hat{\mathbf{z}}$.  The projected covariance is
    compact, giving higher $P_c$.
  \item \textbf{Crossing} ($\mathbf{v}_\text{rel} \approx \hat{\mathbf{n}}$):
    the encounter plane is roughly the R-T plane.  The full 500\,m along-track
    uncertainty \emph{projects into} $\mathbf{C}_\text{2D}$, creating a very
    wide 2D Gaussian.  A 10\,m HBR disk covers a tiny fraction of this
    distribution, giving much lower $P_c$.
\end{itemize}

This geometry effect accounts for 3--6 orders of magnitude difference in $P_c$
between head-on and crossing encounters at the same miss distance (see
\S\ref{sec:findings}).

% ══════════════════════════════════════════════════════════════════════════════
\section{Findings}
\label{sec:findings}
% ══════════════════════════════════════════════════════════════════════════════

\begin{figure}[htb]
  \centering
  \includegraphics[width=\textwidth]{../plots/pc_findings.png}
  \caption{$P_c$ sensitivity study using the Fowler (1993) method.
    In both panels the covariance scale factor is applied equally to \emph{both}
    spacecraft — since the Fowler method sums $\mathbf{C}_1 + \mathbf{C}_2$,
    scaling both by the same factor is equivalent to scaling the combined
    projected covariance by that factor.  This matches the standard assumption
    that both objects have similar tracking quality.
    \textbf{Left:} $P_c$ vs.\ covariance scale for all four scenario types at
    their fixed miss distances.
    Head-on (blue, 200\,m), overtaking (green, 1000\,m), and near-miss (purple,
    10\,m) are all monotone-decreasing: their miss distances are already within
    the spread of the projected Gaussian at default scale, so growing the
    covariance only dilutes it.
    High-$P_c$ crossing (orange, 200\,m) peaks near $3\times$ default before
    falling — it starts in the rising regime where the Gaussian tail barely
    reaches the HBR disk, then crosses into the falling regime (\S\ref{sec:finding1}).
    \textbf{Right:} $P_c$ vs.\ miss distance at fixed default covariance.
    Head-on (blue) stays near $10^{-2}$ out to several kilometers: its encounter
    plane (R-N) excludes the large along-track sigma, keeping the projected
    Gaussian compact.  Crossing (orange) falls much earlier because its encounter
    plane (R-T) includes the full 500\,m along-track sigma, spreading the
    Gaussian far beyond the 10\,m HBR.  Overtaking (green) uses a radial miss
    offset perpendicular to $\mathbf{v}_\text{rel}$; the relevant projected sigma
    is $\sigma_R = 100$\,m, so its curve drops at $\sim$1000\,m — between the
    other two.  Near-miss is omitted: it shares identical crossing geometry and
    its curve lies exactly on the crossing line (\S\ref{sec:finding2}).}
  \label{fig:pc_findings}
\end{figure}

\subsection{Finding 1: $P_c$ vs.\ Covariance Size Is Non-Monotone}
\label{sec:finding1}

Growing the covariance does not always lower $P_c$.  Two competing effects
determine the direction:

\begin{center}
\begin{tabular}{p{3.5cm}p{4cm}p{6cm}}
  \toprule
  Regime & Condition & Effect of growing $\sigma$ \\
  \midrule
  Rising & $d \gg \sigma_\text{2D}$ &
    Gaussian tail barely reaches the HBR disk; growing spreads it toward the
    disk $\Rightarrow$ $P_c$ rises. \\[4pt]
  Falling & $d \ll \sigma_\text{2D}$ &
    Gaussian nearly flat over the disk; growing dilutes it
    $\Rightarrow$ $P_c$ falls. \\
  \bottomrule
\end{tabular}
\end{center}

In all covariance-scale experiments, both spacecraft receive the same scaled
covariance — the scale factor is applied to $\mathbf{C}_1$ and $\mathbf{C}_2$
identically.  Since the Fowler method sums them, this is equivalent to scaling
the combined projected covariance $\mathbf{C}_\text{2D}$ by the same factor.

The high-$P_c$ crossing scenario (204\,m miss, 500\,m/s) is observed to peak
at $\approx$3$\times$ the default covariance.  At tight covariance
($0.1\times$) the 204\,m miss is $>200\,\sigma$ away from the HBR — $P_c$ is
essentially zero.  As the covariance grows, the tail reaches the HBR disk and
$P_c$ rises.  Beyond the peak the Gaussian is so broad it becomes nearly
uniform, and further growth dilutes it.

\textbf{Practical implication:} a new, tighter observation can cause $P_c$ to
\emph{increase} if the operator was previously in the rising regime.  This is
physically correct — better tracking reveals a worse-than-expected situation.

\subsection{Finding 2: Encounter Geometry Determines $P_c$ by Orders of Magnitude}
\label{sec:finding2}

At the same miss distance and the same covariance, head-on conjunctions give
$P_c$ consistently 2--3 orders of magnitude higher than crossing conjunctions.
The difference is purely geometric: what matters is the \emph{direction} of
$\mathbf{v}_\text{rel}$, which controls which axes of the covariance project
into $\mathbf{C}_\text{2D}$.

\begin{itemize}
  \item \textbf{Head-on} ($\mathbf{v}_\text{rel} \approx -\hat{\mathbf{t}}$):
    encounter plane $\approx$R-N.  The large along-track sigma ($\sigma_T = 500$\,m)
    lies along the normal to the plane ($\hat{\mathbf{z}}$) and is excluded
    from the 2D projection entirely.  Only $\sigma_R = 100$\,m and
    $\sigma_N = 50$\,m project in, giving a compact ellipse.
  \item \textbf{Crossing} ($\mathbf{v}_\text{rel} \approx -\hat{\mathbf{n}}$):
    encounter plane $\approx$R-T.  The full $\sigma_T = 500$\,m projects
    directly into $\mathbf{C}_\text{2D}$, creating a wide ellipse.
    A 10\,m HBR disk captures a much smaller fraction of this distribution.
\end{itemize}

The right panel of \Cref{fig:pc_findings} shows this directly: head-on and
crossing both use a 200\,m miss, but head-on's compact projected covariance
keeps $P_c \sim 10^{-3}$ at default scale while crossing's wide projection
gives $P_c \sim 10^{-5}$ — two orders of magnitude lower at the same miss
distance and same covariance magnitude.

\subsection{Finding 3: Near-Miss and Crossing Share Identical Projected Covariances}

The near-miss and crossing scenarios use the same conjunction type (N-dominant
$\mathbf{v}_\text{rel}$, crossing code path) and the same default covariance.
The projected $\mathbf{C}_\text{2D}$ and its eigenvectors are identical — the
$P_c$-vs.-miss curves are indistinguishable and overlap perfectly.  Near-miss
is therefore omitted from the center panel of \Cref{fig:pc_findings} to avoid
redundancy.

What differs is only the designed miss distance used in the test suite: 500\,m
for crossing vs.\ 10\,m for near-miss.  The two scenarios give
$P_c \sim 10^{-14}$ and $P_c \sim 10^{-3}$ respectively — an 11-order-of-magnitude
difference driven entirely by how close the Gaussian center sits to the 10\,m
HBR disk, not by any difference in geometry or covariance.  This is the clearest
demonstration that miss distance and geometry are independent variables in the
Fowler integral.

\subsection{Scenario Comparison}

\begin{figure}[h]
  \centering
  \begin{subfigure}[t]{0.48\textwidth}
    \includegraphics[width=\textwidth]{../plots/comparison_miss_distance.png}
    \caption{Miss distance vs.\ time for all four scenarios.  The TCA is at
             $t = 24$\,h by construction.  Each scenario has a distinct
             approach profile: head-on has the sharpest minimum, overtaking
             has a broad, flat minimum reflecting the low relative speed.}
  \end{subfigure}
  \hfill
  \begin{subfigure}[t]{0.48\textwidth}
    \includegraphics[width=\textwidth]{../plots/comparison_headon_vs_crossing.png}
    \caption{Side-by-side encounter-plane projections for head-on (left) and
             crossing (right) at similar miss distances.  The head-on
             covariance ellipse is compact; the crossing ellipse is highly
             elongated along the T axis, illustrating the geometry effect
             on $P_c$.}
  \end{subfigure}

  \vspace{0.8em}
  \begin{subfigure}[t]{0.48\textwidth}
    \includegraphics[width=\textwidth]{../plots/comparison_crossing_vs_nearmiss.png}
    \caption{Crossing vs.\ near-miss encounter-plane projections.  Both have
             the same orbital geometry (N-dominant $\mathbf{v}_\text{rel}$) but
             different miss distances (500\,m vs.\ 10\,m).  The near-miss
             Gaussian center is close to the HBR disk, giving $P_c \sim 10^{-3}$.}
  \end{subfigure}
  \hfill
  \begin{subfigure}[t]{0.48\textwidth}
    \includegraphics[width=\textwidth]{../plots/comparison_overtaking_2d.png}
    \caption{Overtaking scenario encounter plane.  $\mathbf{v}_\text{rel}$
             is nearly along-track, so the encounter plane is nearly the R-N
             plane — but the miss distance is 1000\,m and $v_\text{rel}$ is
             only 15\,m/s, giving a large projected covariance and low $P_c$.}
  \end{subfigure}
  \caption{Scenario comparison figures from \texttt{plot\_scenario\_comparison.py}.}
  \label{fig:comparison}
\end{figure}

% ══════════════════════════════════════════════════════════════════════════════
\section{Test Suite}
\label{sec:tests}
% ══════════════════════════════════════════════════════════════════════════════

A pytest suite of 70 tests validates all modules.  Session-scoped fixtures
ensure expensive computations (scenario generation, TCA search) run once per
\texttt{pytest} session, not once per test; the full suite completes in
$\approx$35\,s.

\begin{table}[htb]
\centering
\small
\caption{\texttt{test\_conjunction.py} — 26 tests across 5 test classes.}
\label{tab:test_conjunction}
\begin{tabular}{lp{10cm}}
  \toprule
  Test & What it checks \\
  \midrule
  \multicolumn{2}{l}{\textit{Output structure}} \\
  \texttt{test\_returns\_all\_keys} & Dict has exactly the 8 expected keys \\
  \texttt{test\_state\_shapes} & All four ECI state arrays are shape $(6,)$ \\
  \texttt{test\_scalars\_are\_float} & \texttt{miss\_distance} and \texttt{rel\_speed} are Python floats \\
  \texttt{test\_epoch\_ordering} & \texttt{epoch\_start} is strictly before \texttt{epoch\_tca} \\
  \midrule
  \multicolumn{2}{l}{\textit{TCA geometry accuracy}} \\
  \texttt{test\_crossing\_miss\_distance} & Miss distance within 1\,m of requested 500\,m \\
  \texttt{test\_crossing\_rel\_speed} & Relative speed within 0.1\,m/s of requested \\
  \texttt{test\_head\_on\_miss\_distance} & Miss distance within 1\,m of requested 200\,m \\
  \texttt{test\_head\_on\_rel\_speed} & Relative speed within 1\,m/s of requested \\
  \texttt{test\_overtaking\_miss\_distance} & Miss distance within 2\,m of requested 1000\,m \\
  \texttt{test\_near\_miss\_distance} & Miss distance within 0.5\,m of requested 10\,m \\
  \midrule
  \multicolumn{2}{l}{\textit{Physical sanity}} \\
  \texttt{test\_sc1\_altitude\_at\_tca[*]} & SC1 altitude 200--2000\,km (3 scenarios, parametrized) \\
  \texttt{test\_sc1\_speed\_at\_tca[*]} & SC1 orbital speed 6--9\,km/s (3 scenarios, parametrized) \\
  \texttt{test\_two\_spacecraft\_are\_distinct\_at\_t0} & SC1 and SC2 at least 100\,m apart at $T{=}0$ \\
  \midrule
  \multicolumn{2}{l}{\textit{Conjunction type geometry}} \\
  \texttt{test\_crossing\_dominated\_by\_normal} & N component of $\mathbf{r}_\text{rel}$ is largest at TCA \\
  \texttt{test\_head\_on\_dominated\_by\_along\_track} & T component (negative) is largest at TCA \\
  \texttt{test\_overtaking\_dominated\_by\_along\_track} & T component (positive) is largest at TCA \\
  \midrule
  \multicolumn{2}{l}{\textit{Reproducibility}} \\
  \texttt{test\_same\_seed\_same\_result} & Same seed $\Rightarrow$ identical SC1 and SC2 TCA states \\
  \texttt{test\_different\_seeds\_different\_sc2} & SC1 identical; SC2 differs across seeds \\
  \midrule
  \multicolumn{2}{l}{\textit{RTN trajectory sampling}} \\
  \texttt{test\_output\_shape} & Returns array of shape $(n\_\text{samples}, 7)$ \\
  \texttt{test\_time\_column\_starts\_at\_zero} & First time sample is $t=0$ \\
  \texttt{test\_time\_column\_ends\_at\_tca\_hours} & Last time sample matches window length \\
  \texttt{test\_final\_sample\_near\_tca\_miss\_distance} & Last RTN separation matches known miss (within 10\,m) \\
  \bottomrule
\end{tabular}
\end{table}

\begin{table}[htb]
\centering
\small
\caption{\texttt{test\_tca.py} — 13 tests across 4 test classes.}
\label{tab:test_tca}
\begin{tabular}{lp{10cm}}
  \toprule
  Test & What it checks \\
  \midrule
  \multicolumn{2}{l}{\textit{TCA epoch}} \\
  \texttt{test\_crossing\_tca\_timing} & Found TCA within 5\,s of planted TCA \\
  \texttt{test\_head\_on\_tca\_timing} & Same for head-on (15\,km/s, sharp miss profile) \\
  \texttt{test\_overtaking\_tca\_timing} & Same for overtaking (50\,m/s) \\
  \midrule
  \multicolumn{2}{l}{\textit{Miss distance at TCA}} \\
  \texttt{test\_crossing\_miss\_distance} & Found miss within 5\,m of 500\,m \\
  \texttt{test\_head\_on\_miss\_distance} & Found miss within 5\,m of 200\,m \\
  \texttt{test\_overtaking\_miss\_distance} & Found miss within 5\,m of 1000\,m \\
  \texttt{test\_near\_miss\_miss\_distance} & Found miss within 5\,m of 10\,m \\
  \texttt{test\_miss\_distance\_is\_positive} & Miss distance $\ge 0$ \\
  \midrule
  \multicolumn{2}{l}{\textit{States at TCA}} \\
  \texttt{test\_state\_shapes} & Both returned arrays are shape $(6,)$ \\
  \texttt{test\_sc1\_in\_leo} & SC1 altitude 200--2000\,km at found TCA \\
  \texttt{test\_sc1\_speed\_plausible} & SC1 orbital speed in LEO range \\
  \texttt{test\_miss\_distance\_consistent\_with\_states} & State separation matches \texttt{find\_tca} miss (within 1\,m) \\
  \midrule
  \multicolumn{2}{l}{\textit{Edge cases}} \\
  \texttt{test\_coarse\_steps\_parameter} & 50 coarse steps still finds TCA within 30\,s and 10\,m \\
  \bottomrule
\end{tabular}
\end{table}

\begin{table}[htb]
\centering
\small
\caption{\texttt{test\_fowler.py} — 31 tests across 5 test classes.}
\label{tab:test_fowler}
\begin{tabular}{lp{10cm}}
  \toprule
  Test & What it checks \\
  \midrule
  \multicolumn{2}{l}{\textit{Covariance generation}} \\
  \texttt{test\_output\_shapes} & Both covariances are shape $(6, 6)$ \\
  \texttt{test\_symmetric} & Both matrices symmetric to numerical precision \\
  \texttt{test\_positive\_definite} & All eigenvalues positive \\
  \texttt{test\_position\_variance\_order\_of\_magnitude} & ECI diagonal entries in expected range \\
  \texttt{test\_custom\_std} & Trace of position block equals sum of input variances \\
  \midrule
  \multicolumn{2}{l}{\textit{Basic $P_c$ properties}} \\
  \texttt{test\_pc\_in\_unit\_interval} & $P_c \in [0, 1]$ \\
  \texttt{test\_pc\_increases\_as\_miss\_decreases} & Closer miss $\Rightarrow$ higher $P_c$ \\
  \texttt{test\_pc\_increases\_as\_hbr\_increases} & Larger HBR $\Rightarrow$ higher $P_c$ \\
  \texttt{test\_pc\_increases\_as\_covariance\_grows} & Wider covariance $\Rightarrow$ higher $P_c$ (up to saturation) \\
  \texttt{test\_pc\_returns\_float} & Return type is \texttt{float} \\
  \midrule
  \multicolumn{2}{l}{\textit{Limit cases}} \\
  \texttt{test\_zero\_miss\_small\_cov\_near\_one} & Zero miss, HBR $\gg \sigma$ $\Rightarrow$ $P_c > 0.99$ \\
  \texttt{test\_large\_miss\_tiny\_cov\_near\_zero} & Miss $\gg$ HBR, $\sigma \ll$ miss $\Rightarrow$ $P_c < 10^{-10}$ \\
  \texttt{test\_symmetric\_zero\_miss\_depends\_only\_on\_hbr\_and\_sigma} & Different HBR gives different $P_c$ at zero miss \\
  \texttt{test\_large\_covariance\_geometric\_limit} & Larger $\sigma$ lowers $P_c$ in geometric limit \\
  \midrule
  \multicolumn{2}{l}{\textit{Per-scenario magnitudes}} \\
  \texttt{test\_crossing\_pc\_in\_range} & $P_c \in [0, 1]$ for crossing \\
  \texttt{test\_head\_on\_pc\_in\_range} & $P_c \in [0, 1]$ for head-on \\
  \texttt{test\_overtaking\_pc\_in\_range} & $P_c \in [0, 1]$ for overtaking \\
  \texttt{test\_near\_miss\_pc\_higher\_than\_crossing} & 10\,m miss $>$ 500\,m miss \\
  \texttt{test\_pc\_increases\_with\_hbr[*]} & HBR=50\,m gives higher $P_c$ than HBR=5\,m (3 scenarios) \\
  \texttt{test\_slow\_crossing\_pc\_is\_very\_low} & Slow crossing $P_c < 10^{-10}$ \\
  \texttt{test\_high\_pc\_crossing\_magnitude} & Fast crossing $P_c \in [10^{-7}, 10^{-2}]$ \\
  \texttt{test\_head\_on\_pc\_magnitude} & Head-on $P_c \in [10^{-6}, 10^{-1}]$ \\
  \texttt{test\_overtaking\_pc\_magnitude} & Overtaking $P_c \in [10^{-7}, 10^{-1}]$ \\
  \texttt{test\_near\_miss\_pc\_magnitude} & Near-miss $P_c \in [10^{-5}, 10^{-1}]$ \\
  \midrule
  \multicolumn{2}{l}{\textit{Covariance magnitude effects}} \\
  \texttt{test\_head\_on\_pc\_monotone\_decreasing} & Head-on $P_c$ monotone: tight$\to$default$\to$loose$\to$very loose \\
  \texttt{test\_high\_pc\_crossing\_non\_monotone} & Crossing $P_c$ non-monotone: near-zero at tight, peaks near loose \\
  \texttt{test\_near\_miss\_pc\_monotone\_decreasing} & Near-miss $P_c$ monotone across all levels \\
  \texttt{test\_all\_covariance\_levels\_valid} & $P_c \in [0, 1]$ at every covariance level \\
  \midrule
  \multicolumn{2}{l}{\textit{Degenerate inputs}} \\
  \texttt{test\_zero\_relative\_speed\_raises} & \texttt{ValueError} when $|\mathbf{v}_\text{rel}| = 0$ \\
  \bottomrule
\end{tabular}
\end{table}

% ══════════════════════════════════════════════════════════════════════════════
\section{Future Work}
\label{sec:future}
% ══════════════════════════════════════════════════════════════════════════════

The following extensions are planned, in approximate priority order.

\paragraph{Library-grade interfaces.}
The current modules contain test-harness scaffolding (hardcoded constants,
2-spacecraft-at-a-time APIs, implicit unit conventions).  Before contributing
to the \texttt{brahe} library, the modules will be refactored to expose clean,
composable, fully-typed public APIs with NumPy-style docstrings.

\paragraph{Chan (1997) series expansion.}
The Chan series provides a closed-form alternative to the numerical double
integral in \eqref{eq:dblquad}.  It converges in $\sim$10--20 terms and is
$\approx$10$\times$ faster than \texttt{dblquad} for typical $P_c$ values.  It
also handles $P_c < 10^{-10}$ correctly where double-precision quadrature
underflows.  This will be implemented as a drop-in replacement and cross-validated
against the Fowler numerical result.

\paragraph{Monte Carlo baseline.}
A naive Monte Carlo estimator will be added as a validation baseline for all
analytic methods:
\begin{equation}
  \hat{P}_c^\text{MC} = \frac{1}{N}\sum_{i=1}^{N}
    \mathbf{1}\!\left[|\mathbf{x}_i| \le r_\text{hb}\right],
    \qquad \mathbf{x}_i \sim \mathcal{N}(\mathbf{r}_\text{rel},\;\mathbf{C}_\text{pos}).
\end{equation}
$N = 10^6$ samples suffice for $P_c \gtrsim 10^{-4}$.

\paragraph{Importance sampling.}
For $P_c < 10^{-6}$, naive Monte Carlo is impractical ($N \sim 10^8$ required).
Importance sampling with a proposal centered near the HBR boundary will provide
accurate estimates without requiring enormous sample counts.

\paragraph{Patera (2001) line integral.}
An alternative analytic formulation that avoids the 2D Gaussian projection in
favor of a 1D line integral.  Implementing this alongside the Fowler/Chan methods
will enable a structured comparison of the main analytic approaches.

\paragraph{Covariance propagation via STM.}
A state transition matrix (STM) module will propagate the initial covariance
$\mathbf{C}_0$ forward to TCA via finite-difference Jacobian estimation,
replacing the current shortcut of constructing covariances directly at TCA.

% ══════════════════════════════════════════════════════════════════════════════
\section*{References}
% ══════════════════════════════════════════════════════════════════════════════

\begin{enumerate}
  \item Fowler, W.T. (1993). ``A Probability of Collision Analysis for Conjunction
        Assessments.'' AAS Paper 93-177.
  \item Chan, F.K. (1997). ``Spacecraft Collision Probability.'' AAS Paper 97-173.
  \item Patera, R.P. (2001). ``General Method for Calculating Satellite Collision
        Probability.'' \textit{Journal of Guidance, Control, and Dynamics},
        24(4), 716--722.
  \item Alfano, S. (2005). ``A Numerical Implementation of Spherical Object
        Collision Probability.'' \textit{Journal of the Astronautical Sciences},
        53(1), 103--109.
  \item Brahe astrodynamics library. \url{https://github.com/duncaneddy/brahe}
\end{enumerate}

\end{document}
